\documentclass[aspectratio=169,9pt]{beamer}
\mode<presentation> 
\usepackage[utf8]{inputenc}
\usepackage{sansmathaccent} %Für equations
\usepackage{amsmath}
\usepackage{tikz}
\usepackage{multicol}
\usepackage{cancel}
\usepackage{relsize}
\usepackage{graphicx}
\usepackage[table]{xcolor}
\usepackage{fontawesome5}
\usepackage{enumitem}
\usepackage{setspace}
\usepackage[labelfont={bf},font={color=black!40,footnotesize},tablename={Tabelle},figurename={Abbildung}]{caption}
\usetikzlibrary{positioning, arrows.meta, calc, matrix}
\usefonttheme[onlymath]{serif}
\usepackage{booktabs}
%\makeatletter
%\makeatother

\usetikzlibrary{patterns, arrows.meta, calc}
\usepackage{pgfplots}
\pgfplotsset{compat=1.18}
\usetikzlibrary{intersections, pgfplots.fillbetween}\usetikzlibrary{intersections, pgfplots.fillbetween}
\rowcolors{2}{miugray!30}{miugray!10}

\usetheme{Madrid}  %% Themenwahl
\useoutertheme{split} % Alternatively: miniframes, infolines, split
%\useinnertheme{circles}


\setbeamertemplate{itemize item}{\color{miugray}$\blacksquare$}
\setbeamertemplate{itemize subitem}{\color{uniblau}$\bullet$}



%\newcommand{\metermilli}{\mskip3mu \text{mm}}
\newcommand{\cm}{\mskip3mu \text{cm}}
\newcommand{\m}{ \text{m}}
\newcommand{\lite}{ \text{l}}
\newcommand{\kg}{ \text{kg}}
\newcommand{\g}{ \text{g}}
\newcommand{\km}{ \text{km}}
\newcommand{\h}{ \text{h}}
\newcommand{\s}{ \text{s}}
\newcommand{\tonne}{ \text{t}}


\newcommand{\RA}{$\rightarrow$~}
\newcommand{\pp}{\par \vspace*{2mm}}
\newcommand{\ppbig}{\par \vspace*{4mm}}

\definecolor{miudiutex}{HTML}{1d2926}
\definecolor{miugray}{HTML}{b4cccf}
\definecolor{miugray2}{HTML}{4d7365}
\definecolor{white2}{HTML}{f5f7f8}
\definecolor{red}{HTML}{cf1f5c}
\definecolor{green}{HTML}{00A36C}
\definecolor{delphinidin}{HTML}{315794}
\definecolor{uniblau}{HTML}{004e9f}
\definecolor{unigelb}{HTML}{fcba00}
\definecolor{unigrau}{HTML}{909085}
\definecolor{pelargonidin}{HTML}{b33807}
\definecolor{cyanidin}{HTML}{ba0746}
\definecolor{paeonidin}{HTML}{d15ca6}
\definecolor{petunidin}{HTML}{b56fed}
\definecolor{malvidin}{HTML}{8a2d8c}

%\setbeamercolor{palette primary}{bg=pelargonidin!35,fg=white}
%\setbeamercolor{palette secondary}{bg=delphinidin!45,fg=white}
\setbeamercolor{palette primary}{bg=uniblau!70,fg=white}
\setbeamercolor{palette secondary}{bg=uniblau!45,fg=white}
\setbeamercolor{palette tertiary}{bg=unigelb!70,fg=unigrau}
\setbeamercolor{palette quaternary}{bg=miugray,fg=white}
\setbeamercolor{structure}{fg=black} % itemize, enumerate, etc
\setbeamercolor{section in toc}{fg=red} % TOC sections

\setbeamersize{text margin left=1cm,text margin right=1cm}


\useinnertheme{tcolorbox} 
\newcommand{\goodtk}[2]{
\begin{tcolorbox}[
	enhanced,
	colframe=uniblau!75,
	sharp corners,
	boxsep=5pt,
	title=\bfseries \sffamily #1,
	colback=white,
	coltitle=black,
	attach boxed title to top text left={yshift=-\tcboxedtitleheight/2},
	boxed title style={colback=white,colframe=white}
]
	#2
	\end{tcolorbox}
}
	
\newcommand{\antwort}[1]{
\begin{tcolorbox}[
	enhanced,
	colframe=unigrau!65,
	sharp corners,
	boxsep=1pt,
	colback=white,
	coltitle=black,
]
	#1
	\end{tcolorbox}
}


\title{10. Übung}
\author{PD Dr. Helene Loos}
\date{\today}

\begin{document}
%\maketitle
\small



\begin{frame}
\begin{center}

\LARGE Übung 10  \par \normalsize
Prozessbezogene Lebensmitteltechnologie
\end{center}
\small
\end{frame}

















\begin{frame}


\begin{minipage}{.5\textwidth}


{\color{miugray}$\blacksquare ~$} \textbf{Aufgabe 1} \par \vspace*{2mm}

Zwischen dem Serum und den Fettkügelchen von Rohmilch besteht ein Dichteunterschied
von \mbox{$\Delta \rho$ = 99,8 kg/m$^3$.} Die Viskosität der Milch beträgt $\eta$ = 1,79 $\cdot$ 10$^{-3}$ Pa $\cdot$ s. Die Dichte ist mit \mbox{$\rho$ = 1.028 kg/m$^3$} angegeben. 

\begin{enumerate}[label=\alph*)]

\item Berechnen Sie die Aufrahmgeschwindigkeit eines Fettkügelchens mit einem
Durchmesser von \mbox{4 µm} unter der Annahme einer laminaren Strömung.
\item Prüfen Sie mit Hilfe der Reynoldszahl, ob die Annahme einer laminaren Strömung
korrekt war. Erinnerung: Eine laminare Strömung kann angenommen werden, wenn
die Reynoldszahl $<$ 0,5 ist.


\end{enumerate}

\end{minipage}%
\hspace*{4mm}
\begin{minipage}{.5\textwidth}
\antwort{
\begin{align*}
v &= \frac{d^2 \cdot g \cdot (\rho_1 - \rho_2)}{18 \cdot \eta} \\
v &= \frac{(4 \cdot 10^{-6} ~\text{m})^2 \cdot 9,81 ~\text{m/s}^2 \cdot 99,8 ~\text{kg/m}^3}{18 \cdot 1,79 \cdot 10^{-3} ~\text{Pa s}} \\
v &= 4,86 \cdot 10^{-7} ~\text{m/s}
\end{align*}


\begin{align*}
Re &= \frac{\rho \cdot d \cdot v}{\eta} \\
Re &= \frac{1028 ~\text{kg/m}^3 \cdot 4 \cdot 10^{-6} ~\text{m} \cdot 4,86 \cdot 10^{-7} ~\text{m/s}}{1,79 \cdot 10^{-3} ~\text{Pa s}} \\
Re &= 1,116 \cdot 10^{-6}
\end{align*}
\RA Die Annahme einer laminaren Strömung war korrekt

}
\end{minipage}%



\end{frame}










\begin{frame}



\begin{minipage}{.5\textwidth}


{\color{miugray}$\blacksquare ~$} \textbf{Aufgabe 1} \par \vspace*{2mm}

Zwischen dem Serum und den Fettkügelchen von Rohmilch besteht ein Dichteunterschied
von \mbox{$\Delta \rho$ = 99,8 kg/m$^3$.} Die Viskosität der Milch beträgt $\eta$ = 1,79 $\cdot$ 10$^{-3}$ Pa $\cdot$ s. Die Dichte ist mit \mbox{$\rho$ = 1.028 kg/m$^3$} angegeben. 
\begin{enumerate}[label=\alph*)]

\item Die Fettkügelchen (d = 4 µm) haben nach 7.500 Minuten die Oberfläche erreicht. In
welcher Zeit schaffen es die größten Fettkügelchen (d = 5 µm) an die Oberfläche?

\end{enumerate}

\end{minipage}%
\hspace*{4mm}
\begin{minipage}{.5\textwidth}
\antwort{

Strecke berechnen:
\begin{align*}
s = v \cdot t \rightarrow s &= 4,86 \cdot 10^{-7} ~\text{m/s} \cdot 7500 \cdot 60s \\
 &= 0,2187 ~\text{m} \\
\end{align*}
Geschwindigkeit größerer Fettkügelchen berechnen: 
\begin{align*}
v &= \frac{d^2 \cdot g \cdot (\rho_1 - \rho_2)}{18 \cdot \eta} \\
v &= \frac{(5 \cdot 10^{-6} ~\text{m})^2 \cdot 9,81 ~\text{m/s}^2 \cdot 99,8 ~\text{kg/m}^3}{18 \cdot 1,79 \cdot 10^{-3} ~\text{Pa s}} \\
v &= 7,59 \cdot 10^{-7} ~\text{m/s}
\end{align*}
Zeit die größere Fettkügelchen für Strecke $s$ benötigen:
\begin{align*}
t = \frac{s}{v} \rightarrow t &= \frac{0,2187 ~\text{m}}{7,59 \cdot 10^{-7} ~\text{m/s} \cdot 60} \\
&= 4802,4 ~\text{min}
\end{align*}
}
\end{minipage}%


\end{frame}














\begin{frame}
\begin{minipage}{.5\textwidth}


{\color{miugray}$\blacksquare ~$} \textbf{Aufgabe 2} \par \vspace*{2mm}
Aus 70 °C heißer Bierwürze sollen Heißtrub und Hopfentrester in einem runden Absetzbecken
(Füllhöhe 850 mm; Durchmesser 7000 mm) zur Vorklärung absedimentiert werden, bevor das
Gemisch zentrifugiert wird. Die Partikel sind zwischen 30 und 90 µm groß und haben eine
Dichte von 1.050 kg/m$^3$. Für die Dichte und die Viskosität der Würze bei 70 °C nehmen Sie bitte die Kennzahlen für Wasser und laminare Strömungsverhältnisse an.
\pp
Welche Partikel sind nach 40 min vollständig absedimentiert?
\end{minipage}%
\hspace*{4mm}
\begin{minipage}{.5\textwidth}
\antwort{

Geschwindigkeit die Partikel brauchen um nach 40 min vollständig absedimentiert zu sein:
\begin{align*}
v &= 0,85~\text{m} / ~2400~\text{s} \\
&= 3,5 \cdot 10^{-4} ~\text{m/s}
\end{align*}
Stokes umstellen:
\begin{align*}
v &= \frac{d^2 \cdot g \cdot (\rho_1 - \rho_2)}{18 \cdot \eta} \\
d &= \sqrt{\frac{v \cdot 18 \cdot \eta}{g \cdot \Delta p}} \\
d &= \sqrt{\frac{3,5 \cdot 10^{-4} ~\text{m/s} \cdot 18 \cdot 1 \cdot 10^{-3} ~\text{Pa s}}{9,81 ~\text{m/s} \cdot 50 ~\text{kg/m}^3}} \\
% &= 7,2 \cdot 10^{-5} ~\text{m} \\
% &= 72 ~\text{µm}
&= 1,13 \cdot 10^{-4} ~\text{m} \\
&= 113 ~\text{µm}
\end{align*}
Alle Partikel die größer als 113 µm sind, sedimentieren vollständig.
}
\end{minipage}%
\end{frame}









\begin{frame}


\begin{minipage}{.5\textwidth}


{\color{miugray}$\blacksquare ~$} \textbf{Aufgabe 3} \par \vspace*{2mm}
Milch soll mittels Mikrofiltration entkeimt werden. Hierfür wählen Sie ein geeignetes
Filterelement mit einem Porendurchmesser von 0,2 µm, einer Länge von 20 cm und einer
Porenanzahl von 1.000.000. Gehen Sie dabei bitte von einer laminaren Strömung aus. Die
Viskosität der Milch beträgt 2 mPa$\cdot$s. Durch Anlegen eines Vakuums wird eine Druckdifferenz
von 0,5 bar erzeugt. Wie groß ist die Volumenstromstärke der Milch durch das Filterelement? 
\end{minipage}%
\hspace*{4mm}
\begin{minipage}{.5\textwidth}
\antwort{
\begin{align*}
\dot{V} &= \frac{\pi \cdot \Delta p \cdot R^4}{8 \cdot l \cdot \eta} \\
\dot{V} &= \frac{\pi \cdot 5 \cdot 10^4 ~\text{Pa} \cdot (1 \cdot 10^{-7} ~\text{m})^4}{8 \cdot 0,2 ~\text{m} \cdot 2 \cdot 10^{-3} ~\text{Pa s}} \\
&= 4,91 \cdot 10^{-21} ~\text{m}^3/\text{s}
\end{align*}
Dadurch, dass wir aber eine Anzahl von 1.000.000 Poren haben, müssen wir den Volumenstrom $\dot{V}$ nochmal mit der Anzahl multiplizieren:
\begin{align*}
4,91 \cdot 10^{-21} ~\text{m}^3/\text{s} \cdot 1.000.000 = 4,91 \cdot 10^{-15} ~\text{m}^3/\text{s}
\end{align*}

}
\end{minipage}%

\end{frame}









\begin{frame}
\begin{minipage}{.5\textwidth}


{\color{miugray}$\blacksquare ~$} \textbf{Aufgabe 4} \par \vspace*{2mm}
Aus einer wässrigen Suspension sollen Sie mittels einer Dekanterzentrifuge (Durchmesser $d$ = 500 mm, Länge $l$ = 450 mm) den Feststoff zurückgewinnen. Die Zentrifuge wird mit 40 \% Füllung betrieben.  
\begin{itemize}
\item Berechnen Sie den mittleren Durchmesser des sich in der Zentrifuge einstellenden Flüssigkeitsringes.
\item Was ist der Vorteil einer Dekanterzentrifuge gegenüber anderen Zentrifugentypen?
\end{itemize}
\end{minipage}%
\hspace*{4mm}
\begin{minipage}{.5\textwidth}
\antwort{
40 \% des Gesamtvolumens entsprechen dem Flüssigkeitsring:
\begin{align*}
\pi \cdot (r_a^2-r_i^2) \cdot h = 0,4 \cdot \pi \cdot r^2 \cdot h \\
 r_a^2-r_i^2 &= 0,4 r^2 \\ 
 0,6 \cdot r_a^2 &= r_i^2 \\ 
 \sqrt{0,6} \cdot r_a &= r_i
\end{align*}
Werte einsetzen: 
\begin{align*}
r_i &= \sqrt{0,6} \cdot 0,25~\text{m} \\
&= 0,193 ~\text{m}
\end{align*}
Mittlerer Durchmesser:
\begin{align*}
\left( \frac{0,25~\text{m}+0,193~\text{m}}{2} \right) \cdot 2 = 0,443~\text{m}
\end{align*}

%\pp 
%
%\textbf{Vorteile Dekanterzentrifugen} \par 
%Dekanterzentrifugen arbeiten kontinuierlich und können große Feststoffmengen tolerieren. So können Teller - und Röhrenzentrifugen beispielsweise nur kleinste Mengen Feststoff vertragen. Wobei hierbei auch die Tellerzentrifuge nochmal ein bisschen mehr Feststoff toleriert als die Röhrenzentrifuge. 
}
\end{minipage}%
\end{frame}



\begin{frame}

\begin{table}
\caption{Eigenschaften von Zentrifugen nach Brennan}
\begin{spacing}{2}

\begin{tabular}{llll} 
\toprule
\textbf{Zentrifuge} & \textbf{Feststoffanteil} & \textbf{Kontinuierlich} & \textbf{Filterung feiner Partikel} \\ \hline
\multicolumn{4}{l}{\textit{Flüssig-Flüssig Zentrifugen}} \\
Röhrenzentrifuge & $<$ 0,5 \% & Bedingt & Ja \\
Tellerzentrifuge & $<$ 1 - 2,5 \% & Bedingt & Ja \\ 
\hline
\multicolumn{4}{l}{\textit{Flüssig-Fest Zentrifugen}} \\
Dekanterzentrifuge & bis zu 90\% & Ja & Nein (nur Partikel $>$ 2 µm) \\ \bottomrule
\end{tabular}

\end{spacing}
\end{table}
\end{frame}






\end{document}
